\section{When the answer is no}
\label{sec:ch05-when-the-answer-is-no}
\index{errors!field errors|(}

GraphQL does not have status codes. A request that reaches the execution engine
answers 200, and that is true of a request where every field resolved and of a
request where none of them did. Whatever a client learns about failure, it
learns from the body.

The body has two places to put it, and choosing between them is most of what
error design is. This section is the first place: the \code{errors} array,
which every GraphQL server has and which most services over-use.

\subsection*{An error is per field, not per request}

\index{errors!non-terminating}%
The specification calls these field errors, and the important word in the
definition is non-terminating: a resolver that fails returns null for its own
field, an entry is added to \code{errors}, and every other field in the query
goes on resolving~\autocite{chillicream2026errors}. Ask Mosaic's \code{node}
field for an identifier naming a type nothing is registered for:

\begin{minted}[fontsize=\footnotesize]{json}
{"errors":[{"message":"There is no node resolver registered for type `Widget`.",
  "path":["node"]}],
 "data":{"node":null}}
\end{minted}

\code{data} is there. The field that failed is null and everything beside it
would have resolved. A client that reads the presence of \code{errors} as total
failure is throwing away answers it was given.

\index{errors!null propagation}%
Now the same failure where the field cannot be null. The mutation that submits a
review takes an argument typed to accept only a product, and this hands it a
customer:

\begin{minted}[fontsize=\footnotesize]{json}
{"errors":[{"message":"The node id type name `Customer` does not match
  the expected type name `Product`.",
  "path":["submitReview"]}],
 "data":null}
\end{minted}

The field returns \code{SubmitReviewPayload!}, so there is nowhere to write a
null, and the error walks up until it finds somewhere: the root. Same class of
failure, same array, and the entire difference in blast radius is one
exclamation mark. Chapter~\ref{ch:data-without-the-n-plus-1} met this from the
other direction, when a missing \code{totalCount} took a whole page of products
with it. Nullability is a statement about how far a single failure is allowed
to travel, and it is the only lever you have over that distance.

\subsection*{What a client is told by default}

\index{errors!masking}%
Both of those are HotChocolate's own errors, which is why they say something
useful. An exception escaping one of your resolvers is treated quite
differently: the documentation's example of it, since nothing in Mosaic throws
an unhandled exception on any input I could find, is a response whose message
reads \enquote{Unexpected Execution Error} and nothing
else~\autocite{chillicream2026errors}.

The message is gone. That default is correct and
it is worth understanding rather than switching off, because the first thing
most people do when they meet it is switch it off. An exception message is
written for whoever reads the log, and it routinely contains a connection
string, a file path, an internal identifier or a stack of type names that
describes the shape of your system to somebody who was guessing. The engine
cannot tell which exceptions were written for a client, so it assumes none of
them were.

\index{errors!ErrorBuilder}%
Two supported ways to say more, neither of which Mosaic needs yet and both of
which the documentation is the authority on. An error filter catches a typed
exception on its way out and rewrites it:

\begin{minted}[fontsize=\footnotesize]{csharp}
builder
    .AddGraphQL()
    .AddErrorFilter(error =>
    {
        if (error.Exception is NotFoundException ex)
        {
            return ErrorBuilder
                .FromError(error)
                .SetMessage(ex.Message)
                .SetCode("NOT_FOUND")
                .Build();
        }

        return error;
    });
\end{minted}

Or a resolver throws \code{GraphQLException}, which is passed through as
written~\autocite{chillicream2026errors}. The difference between the two is
where the decision lives. A filter is central and catches a category of failure
across the whole schema; an exception is local and is the resolver saying this
particular message is for you. Errors are immutable, and the methods that look
like mutations each return a new instance, which is why
\code{ErrorBuilder.FromError} is the documented way to change more than one
thing at once~\autocite{chillicream2026errors}.

\code{SetCode} is the part to copy whichever route you take. A message is
written for a person, gets reworded, translated and shortened, and any client
that branched on its text breaks the day somebody improves it. A code is written
for a program.

\subsection*{The array is the wrong place for most of what goes wrong}

\index{errors!when to use the array}%
Here is the position this book takes, and it is a position rather than a
consensus.

The \code{errors} array is for failures the client cannot do anything about: a
database that was unreachable when the resolver asked, or a bug that threw.
What those have in common is that there is no branch for a client to take and
nothing to say to the person at the keyboard beyond the fact that something went
wrong.

\enquote{You have already reviewed this product} is not that. It is the correct,
expected outcome of a legitimate request, the client knows exactly what to do
with it, and burying it in a channel meant for stack traces means the client has
to string-match its way back out.

Two things go wrong when a domain failure is put in the array, and the second is
the one that compounds. The type system never sees it, so no client generator
knows the failure can happen and no compiler will tell a client it has not
handled it. And the shape of the error is then a convention rather than a
contract: adding a field to it, or renaming one, is a change no schema check can
see, which is the same blindness
section~\ref{sec:ch05-one-name-for-one-thing} ran into with the identifier
format.

\index{errors!two channels}%
So Mosaic ends up with two channels, and the line between them is about whether
the caller has a decision to make. A failure with a branch behind it is data and
belongs in the response. A failure with no branch behind it is an error and
belongs in the array.

That distinction needs a mutation to demonstrate, and Mosaic has never had one.
\index{errors!field errors|)}
